% CẤU HÌNH CHUNG - theo Phụ lục 3, QĐ 449/QĐ-HV (PTIT). Biên dịch bằng XeLaTeX
\usepackage{fix-cm}
\usepackage{fontspec}
\IfFontExistsTF{Times New Roman}
  {\setmainfont{Times New Roman}}
  {\setmainfont{TeX Gyre Termes}}   % thay thế khi máy không có Times New Roman
% Không dùng babel-vietnamese (nhiều bản TeX không có file .ldf): đặt tên nhãn thủ công,
% tắt ngắt từ vì tiếng Việt không cần gạch nối cuối dòng.
\hyphenpenalty=10000 \exhyphenpenalty=10000 \tolerance=2000
\renewcommand{\figurename}{Hình}\renewcommand{\tablename}{Bảng}
\usepackage{amsmath}
\usepackage{exscale}
\DeclareFontShape{OMX}{cmex}{m}{n}{%
  <-8>sfixed * cmex7
  <8> cmex8
  <9> cmex9
  <10> <10.95> <12> <13> <14.4> <17.28> <20.74> <24.88>cmex10
}{}
\usepackage{graphicx,xcolor,array,booktabs,longtable,multirow,enumitem}
\usepackage{tabularx,fancyvrb}
\IfFileExists{ragged2e.sty}{\usepackage{ragged2e}}{\newcommand{\RaggedRight}{\raggedright}}
\usepackage{float,tikz}
\usetikzlibrary{positioning,arrows.meta,calc,shapes.geometric,shapes.multipart,fit,backgrounds}

% --- Khổ giấy: A4; top/bottom/right 2 cm, left 3 cm; header/footer ~1 cm
\usepackage[a4paper,top=2cm,bottom=2cm,left=3cm,right=2cm,
            headheight=17pt,headsep=1cm,footskip=1cm]{geometry}

% --- Cỡ chữ 13 pt, giãn dòng 1.2, before/after = 3 pt
\makeatletter
\renewcommand\normalsize{\@setfontsize\normalsize{13}{15.6}%
  \abovedisplayskip 8\p@ \@plus2\p@ \@minus4\p@
  \abovedisplayshortskip \z@ \@plus3\p@
  \belowdisplayshortskip 6\p@ \@plus3\p@ \@minus3\p@
  \belowdisplayskip \abovedisplayskip
  \let\@listi\@listI}
\renewcommand\small{\@setfontsize\small{12}{14.4}}
\renewcommand\footnotesize{\@setfontsize\footnotesize{11}{13.2}}
\renewcommand\scriptsize{\@setfontsize\scriptsize{9.5}{11.4}}
\makeatother
\normalsize
\linespread{1}
\setlength{\parindent}{0pt}
\setlength{\parskip}{6pt}          % 3 pt trước + 3 pt sau
\sloppy
\setlist{nosep,itemsep=3pt,topsep=3pt}

% --- Đánh số: 1.1, 1.1.1, 1.1.1.a ; công thức theo chương
\renewcommand{\thesubsubsection}{\thesubsection.\alph{subsubsection}}
\setcounter{secnumdepth}{3}
\setcounter{tocdepth}{3}
\numberwithin{equation}{chapter}

% --- Tiêu đề chương / mục
\usepackage{titlesec}
\titleformat{\chapter}[block]{\normalfont\bfseries\fontsize{15}{18}\selectfont\centering}
  {CHƯƠNG \thechapter:}{0.4em}{}
\titleformat{name=\chapter,numberless}[block]{\normalfont\bfseries\fontsize{15}{18}\selectfont\centering}{}{0pt}{}
\titlespacing*{\chapter}{0pt}{0pt}{18pt}
\titleformat{\section}{\normalfont\bfseries\normalsize}{\thesection}{0.6em}{}
\titleformat{\subsection}{\normalfont\bfseries\normalsize}{\thesubsection}{0.6em}{}
\titleformat{\subsubsection}{\normalfont\bfseries\itshape\normalsize}{\thesubsubsection}{0.6em}{}
\titlespacing*{\section}{0pt}{12pt}{4pt}
\titlespacing*{\subsection}{0pt}{9pt}{3pt}
\titlespacing*{\subsubsection}{0pt}{6pt}{3pt}

% --- Header / Footer
% Header trái: "Đồ án tốt nghiệp Đại học"; header phải: chương và tên chương.
% Footer trái: tên sinh viên, tên lớp; footer phải: số trang.
\usepackage{fancyhdr}
\pagestyle{fancy}                  % để fancyhdr khởi tạo xong rồi mới định nghĩa lại các mark bên dưới
\newcommand{\DATNstyle}{%
  \fancyhf{}
  \fancyhead[L]{\small Đồ án tốt nghiệp Đại học}
  \fancyhead[R]{\small\rightmark}
  \fancyfoot[L]{\small\FooterSV}
  \fancyfoot[R]{\small\thepage}
  \renewcommand{\headrulewidth}{0.4pt}
  \renewcommand{\footrulewidth}{0.4pt}}
\fancypagestyle{datn}{\DATNstyle}
\fancypagestyle{plain}{\DATNstyle}          % trang đầu chương cũng có header/footer
\fancypagestyle{nonum}{\DATNstyle\fancyfoot[R]{}}   % trang nhiệm vụ: không số trang

% Chương không đánh số (LỜI CẢM ƠN, MỞ ĐẦU, ...): ghi tiêu đề IN HOA
\newcommand{\unnumchapter}[1]{%
  \chapter*{#1}\addcontentsline{toc}{chapter}{#1}\markboth{}{#1}}
% Chương có đánh số: \chuong{TIÊU ĐỀ IN HOA}{Tiêu đề viết thường cho header}
% (header luôn ghi "Chương n: tên chương"; tên mục không ghi đè lên header)
\newcommand{\TenChuongHeader}{}
\newcommand{\chuong}[2]{\renewcommand{\TenChuongHeader}{#2}\chapter{#1}}
\renewcommand{\chaptermark}[1]{\markboth{}{Chương \thechapter: \TenChuongHeader}}
\renewcommand{\sectionmark}[1]{}
% Phụ lục: không đánh số chương, bảng/hình/sơ đồ đánh số PL.1, PL.2, ...
\newcommand{\phuluc}{%
  \unnumchapter{PHỤ LỤC}%
  \setcounter{table}{0}\setcounter{figure}{0}\setcounter{diagram}{0}%
  \renewcommand{\thetable}{PL.\arabic{table}}%
  \renewcommand{\thefigure}{PL.\arabic{figure}}%
  \renewcommand{\thediagram}{PL.\arabic{diagram}}}
\makeatletter
\newcommand{\mucphuluc}[1]{\par\addvspace{9pt}%
  {\noindent\bfseries #1\addcontentsline{toc}{section}{#1}\par}\nobreak\vskip 2pt\@afterheading}
\makeatother

% --- Hình, bảng, sơ đồ: chú thích dưới hình/sơ đồ, trên bảng
% Bảng ngắn, hình và sơ đồ dùng [!htbp] (thả nổi) để không để trống cuối trang; bảng dài dùng longtable.
% flafter: không đặt hình/bảng trước chỗ nó được khai báo; \FloatBarrier (placeins): chặn không cho trôi qua một điểm.
\usepackage{flafter}
\IfFileExists{placeins.sty}{\usepackage{placeins}}{\providecommand{\FloatBarrier}{}}
\renewcommand{\topfraction}{0.9}
\renewcommand{\bottomfraction}{0.8}
\renewcommand{\textfraction}{0.08}
\renewcommand{\floatpagefraction}{0.8}
\setcounter{topnumber}{3}\setcounter{bottomnumber}{2}\setcounter{totalnumber}{4}
\setlength{\textfloatsep}{12pt plus 3pt minus 2pt}
\setlength{\intextsep}{10pt plus 3pt minus 2pt}
\setlength{\floatsep}{10pt plus 3pt minus 2pt}
\makeatletter   % trang chỉ chứa hình/bảng: xếp từ đầu trang thay vì canh giữa
\setlength{\@fptop}{0pt}\setlength{\@fpsep}{14pt plus 2pt}\setlength{\@fpbot}{0pt plus 1fil}
\makeatother
\usepackage{caption}
\captionsetup{font=small,labelfont=bf,labelsep=colon,skip=6pt}
\captionsetup[table]{position=top}
\captionsetup[longtable]{position=top}
\newfloat{diagram}{htbp}{lod}[chapter]
\floatname{diagram}{Sơ đồ}
\captionsetup[diagram]{position=bottom}
% Ghi chú nguồn dưới mỗi hình/bảng/sơ đồ (theo Mẫu 7)
\newcommand{\nguon}[1]{\par\vspace{2pt}{\centering\small\itshape Nguồn: #1\par}}

% --- Bảng: kiểu cột canh trái có độ rộng cố định, chữ nhỏ trong bảng
\newcolumntype{L}[1]{>{\RaggedRight\arraybackslash}p{#1}}
\newcolumntype{C}[1]{>{\centering\arraybackslash}p{#1}}
\newcolumntype{Y}{>{\RaggedRight\arraybackslash}X}
\newcommand{\bang}{\small\renewcommand{\arraystretch}{1.18}}
\newcommand{\td}[1]{\textbf{#1}}     % ô tiêu đề cột
\setlength{\LTpre}{6pt}\setlength{\LTpost}{2pt}

% --- Mã nguồn, định danh kỹ thuật
\newcommand{\ma}{\path}              % \ma{/api/v1/agent/command}: chữ đánh máy, tự xuống dòng tại / . _
\DefineVerbatimEnvironment{code}{Verbatim}
  {fontsize=\scriptsize,frame=single,framesep=4pt,rulecolor=\color{black!45},
   xleftmargin=1pt,xrightmargin=1pt,baselinestretch=0.92}
\fvset{listparameters=\setlength{\topsep}{4pt}\setlength{\partopsep}{0pt}}

% --- Ảnh chụp màn hình: \anh[độ rộng]{tên tệp trong figures/}
\newcommand{\anh}[2][0.86\textwidth]{{%
  \setlength{\fboxsep}{0pt}\setlength{\fboxrule}{0.3pt}%
  \IfFileExists{figures/#2}
    {\fbox{\includegraphics[width=#1]{figures/#2}}}
    {\fbox{\parbox[c][5cm][c]{#1}{\centering\small Chèn ảnh: \texttt{figures/\detokenize{#2}}}}}}}

% --- Kiểu vẽ sơ đồ TikZ dùng chung
\tikzset{
  hop/.style={draw,rounded corners=2pt,minimum height=0.8cm,align=center,font=\footnotesize,inner sep=4pt},
  hopdam/.style={hop,fill=black!8},
  nhom/.style={draw,dashed,rounded corners=3pt,inner sep=7pt},
  mui/.style={-{Stealth[length=2.2mm]},line width=0.5pt},
  muihai/.style={{Stealth[length=2.2mm]}-{Stealth[length=2.2mm]},line width=0.5pt},
  nhan/.style={font=\scriptsize,align=center,inner sep=1.5pt},
  nhannen/.style={nhan,fill=white},
  seqhead/.style={draw,fill=black!8,minimum height=0.95cm,align=center,font=\scriptsize\bfseries,inner sep=3pt},
  seqline/.style={dashed,line width=0.4pt,black!55},
  seqmsg/.style={-{Stealth[length=2mm]},line width=0.5pt},
  seqret/.style={-{Stealth[length=2mm]},dashed,line width=0.5pt},
  seqtxt/.style={font=\scriptsize,align=center,inner sep=1.5pt,fill=white},
  uc/.style={draw,ellipse,align=center,font=\footnotesize,inner sep=1.5pt,minimum height=0.85cm,minimum width=3.3cm},
  thucthe/.style={draw,rectangle split,rectangle split parts=2,rectangle split part fill={black!8,white},
                  align=left,font=\scriptsize,inner sep=3pt,text width=3.7cm},
  trangthai/.style={draw,rounded corners=8pt,minimum height=0.95cm,minimum width=2.6cm,align=center,font=\footnotesize},
}
% Thông điệp trong sơ đồ tuần tự: \tn{y}{x đầu}{x cuối}{nội dung} (nét liền), \tnve (nét đứt, trả về),
% \tntu[neo]{y}{x}{nội dung} (bước tự xử lý, vẽ thành hộp trên đường sống)
\newcommand{\tn}[4]{\draw[seqmsg] (#2,-#1) -- node[seqtxt,above=0.6pt]{#4} (#3,-#1);}
\newcommand{\tnve}[4]{\draw[seqret] (#2,-#1) -- node[seqtxt,above=0.6pt]{#4} (#3,-#1);}
\newcommand{\tntu}[4][center]{\node[seqtxt,draw,rounded corners=1pt,inner sep=2.5pt,anchor=#1] at (#3,-#2) {#4};}

% --- Mục lục & Danh mục bảng, sơ đồ, hình vẽ
\usepackage{tocloft}
\renewcommand{\cftchapfont}{\small\bfseries}
\renewcommand{\cftchappagefont}{\small\bfseries}
\renewcommand{\cftchappresnum}{CHƯƠNG~}
\renewcommand{\cftchapaftersnum}{:}
\setlength{\cftchapnumwidth}{5.6em}
\renewcommand{\cftchapleader}{\cftdotfill{\cftdotsep}}
\renewcommand{\cftdotsep}{1.5}
\setlength{\cftbeforechapskip}{4pt}
\renewcommand{\cftsecfont}{\small\bfseries}
\renewcommand{\cftsecpagefont}{\small}
\renewcommand{\cftsubsecfont}{\small}
\renewcommand{\cftsubsubsecfont}{\small}
\cftsetindents{section}{1.5em}{2.5em}
\cftsetindents{subsection}{4em}{3.2em}
\cftsetindents{subsubsection}{7.2em}{3.8em}
\setlength{\cftbeforesecskip}{2pt}
\setlength{\cftbeforesubsecskip}{1pt}
\renewcommand{\cfttabpresnum}{BẢNG~}
\renewcommand{\cftfigpresnum}{HÌNH~}
\newcounter{loddepth}\setcounter{loddepth}{1}
\newlistentry{sodo}{lod}{0}          % kiểu dòng mục lục cho "Sơ đồ"
\makeatletter\let\l@diagram\l@sodo\makeatother
\renewcommand{\cftsodopresnum}{SƠ ĐỒ~}
\setlength{\cftsodoindent}{1.5em}
\setlength{\cfttabnumwidth}{6.4em}
\setlength{\cftfignumwidth}{6.4em}
\setlength{\cftsodonumwidth}{6.4em}
\renewcommand{\cfttabfont}{\small}
\renewcommand{\cftfigfont}{\small}
\renewcommand{\cftsodofont}{\small}
\setlength{\cftbeforetabskip}{2pt}
\setlength{\cftbeforefigskip}{2pt}
\setlength{\cftbeforesodoskip}{2pt}
\setlength{\cftparskip}{0pt}

% --- Siêu liên kết (không tô màu để in)
\usepackage[hidelinks,unicode,hypertexnames=false,pdftitle={\TenDeTai}]{hyperref}
\Urlmuskip=0mu
